\bindcontent{
\section*{Question 1}
A transportation researcher at UCF wants to study whether a new carpooling initiative reduces students' average weekly commute time.
The true average commute time of all UCF students is $22$ minutes, but this value is unknown to the researcher.
She selects $150$ UCF students who volunteer to participate in the initiative. After $6$ weeks, the average commute time of the sample decreases from $24$ minutes to $18$ minutes.

\begin{enumerate}[label=(\alph*)]
    \item Identify the following:
    \begin{enumerate}[label=\roman*.]
        \item Population
        \item Sample
        \item Parameter
        \item Statistic
    \end{enumerate}

    \item Classify each of the following variables as \textbf{qualitative} or \textbf{quantitative}. If quantitative, state whether it is \textbf{discrete} or \textbf{continuous}.
    \begin{enumerate}[label=\roman*.]
        \item Weekly commute time (in minutes)
        \item Primary mode of transportation (e.g., Car, Bus, Walk)
        \item Number of days a student commutes to campus per week
        \item Whether a student uses the carpooling initiative (Yes/No)
    \end{enumerate}

    \item Identify the level of measurement for each variable:
    \begin{enumerate}[label=\roman*.]
        \item Weekly commute time
        \item Primary mode of transportation
        \item Number of days a student commutes
    \end{enumerate}

    \item Is this study an \textbf{observational study} or a \textbf{designed experiment}? Justify your answer.

    \item Can the researcher conclude that the carpooling initiative \textbf{caused} the reduction in commute times? Explain clearly.

    \item Suppose instead the researcher randomly selected students from all UCF students using a student ID list.
    \begin{enumerate}[label=\roman*.]
        \item What type of sampling method is this?
        \item Why is this method preferred over the original method used?
    \end{enumerate}
\end{enumerate}
}
